\section{A fresh adversary search against the frozen configuration}
\label{sec:exploit}

The protocol's fifth failure condition is that any fitted adversary clears the detection floor against
the frozen configuration with a replicated sign: ``then v0.7 closed nothing and opened something, and
the exploit is the headline''. Eight independent searches were run to find out.

\begin{table}[t]\centering
\caption{Eight independent adversary searches against \frozen{}. A positive edge would mean the
adversary beats the target; every one is negative. Each arm was fitted on a stream disjoint from the
banks it is evaluated on.}
\label{tab:adversaries}
{\small\begin{tabularx}{\linewidth}{@{}l L l r c r c@{}}
\toprule
ID & Class & Objective & Edge (pp) & 95\% CI & Per bank & Repl. \\
\midrule
Z01 & C1 in-class & win & $-2.91$ & $[-3.55, -2.27]$ & $-2.76$ / $-3.06$ & yes \\
Z02 & C2 extended & win & $-3.90$ & $[-4.52, -3.28]$ & $-3.80$ / $-4.00$ & yes \\
Z03 & C1 & declerr & $-10.70$ & $[-11.32, -10.09]$ & $-10.65$ / $-10.76$ & yes \\
Z04 & C1 & events & $-9.64$ & $[-10.27, -9.01]$ & $-9.83$ / $-9.46$ & yes \\
Z05 & C1 & forced & $-7.14$ & $[-7.76, -6.52]$ & $-7.17$ / $-7.11$ & yes \\
Z06 & C2 wide sigma & win & $-5.25$ & $[-5.88, -4.62]$ & $-4.93$ / $-5.57$ & yes \\
Z07 & C1 v05 basin & win & $-5.21$ & $[-5.84, -4.58]$ & $-4.85$ / $-5.58$ & yes \\
Z08 & C5 white-box & win & $-3.56$ & $[-4.19, -2.92]$ & $-3.95$ / $-3.17$ & yes \\
\bottomrule
\end{tabularx}
}
\end{table}

\vsevenBfourFits{} of \vsevenBfourArms{} searches produced a weight vector and all
\vsevenBfourArms{} were evaluated. \textbf{Not one arm beats the frozen configuration.} Every one
loses --- by between \vsevenBfourBestLoss{} and \vsevenBfourWorstLoss{} points --- every sign
replicates, and the largest upper bound anywhere is \vsevenBfourWorstUpper{}, which is
\vsevenBfourFloorGap{} points below the floor. The preregistered claim that no fresh adversary
exploits the configuration passes.

\paragraph{The pass condition and the fail condition are not complements, and both were computed.}
``Every arm's upper bound below its floor'' and ``any arm clears its floor with a replicated sign''
leave a band in which the protocol determines no verdict. Had that band been occupied the protocol
required the evaluation to stop rather than to resolve it by choosing a reading. The band is
\emph{empty} --- every upper bound is below its floor --- so the pass condition is met outright and no
reading had to be chosen. This is recorded because it was a live risk, not because it materialised.

\paragraph{Four objectives, not one.} The searches maximise win rate, declaration error, event count
and forced-endgame incidence. The declaration-objective arm is read against the declaration-family
floor of \vsevenDeclFloor{} rather than \vsevenPrimaryFloor{}, because that is its whole objective; its
upper bound is far enough below both that the distinction changes no verdict here. Two arms
deliberately depart from the template --- one with a wide search variance, one seeded in the \vfive{}
basin rather than the incumbent's --- so that the eight are not eight runs of one search.

\subsection{The white-box arm, and its two limits}
\label{sec:exploit-whitebox}

\texttt{Zeight} is the class the threat model of \S\ref{sec:threat} was written for: an inverter that
reads the target's own deterministic policy and reconstructs the deal posterior from the public
transcript. It reaches \vsevenZeightEdge{} $[\vsevenZeightLo, \vsevenZeightHi]$ --- a loss, replicated.

Two limits are stated because this is the cell a reader will weigh most heavily.

First, \textbf{\vsevenZeightInert{} of its fitted coordinates are inert.} The white-box base derives from the
\vsix{} agent and reads only that policy's \vsevenVsixCoords{} coordinates, while the search extends its
vector to \vsevenCoordinates{} for any base of this family. The search is therefore over \vsevenZeightLive{}
live coordinates plus the inverter the base fixes, and a reader should price it as that rather than as
a \vsevenCoordinates{}-dimensional search.

Second, \textbf{the budget is small}: \vsevenBfourGens{} generations of a
\vsevenBfourPop{}-member population over \vsevenBfourDeals{} deals a candidate. \texttt{Zeight}
is evidence that this white-box inverter, at this budget, does not exploit the frozen configuration.
It is not evidence that no white-box attack can. \S\ref{sec:limitations-lbr} says so again, because
this is the exact place the corpus has previously turned a null result into a win.

\subsection{The winner's curse, visible in the fitting-to-holdout gap}
\label{sec:exploit-curse}

The gap between fitting performance and holdout performance is large and in the expected direction:
the in-class arm reached \vsevenZoneFitMargin{}~points of paired margin in its final generation on
its own fitting stream --- \vsevenZoneFitPeak{} at its best generation --- and lands at
\vsevenZoneEdge{} on disjoint banks. That is the arithmetic the disjointness rule exists to expose,
and it is why an exploitability number measured on its own fitting stream is not a measurement.

A second sample of the same finding sits in \S\ref{sec:results-controls}: the responders fitted for
the planted-edge controls, on a different bank and against a handicapped target, also lose to
unhandicapped \frozen{} in every row. They maximise win rate, as four of these eight arms do, so it
is not an independent \emph{method}.
